\subsection{Secuencia mágica}
\label{sec:secuencia}

\subsubsection{Descripción del problema}

Una secuencia mágica de longitud $n$ es una tupla $x_0, \dots, x_{n-1}$ donde cada componente $x_i$ cuantifica el número de apariciones del valor $i$ en toda la secuencia. El enunciado pide un modelo parametrizado en $n$ que enumere todas las secuencias mágicas de longitud $n$, y luego agregarle dos restricciones redundantes dadas.

\subsubsection{Modelo}
Formalmente se modela como el CSP $\mathcal{P} = \langle \mathcal{V}, \mathcal{D}, \mathcal{C} \rangle$ \parencite{apt2003principles}:

\paragraph{Parámetros}
\begin{itemize}
    \item $n \in \mathbb{N}$: longitud de la secuencia ($n \ge 1$).
\end{itemize}

\paragraph{Variables y dominios}
Se define el conjunto de variables de decisión:
\[
\mathcal{V} = \{x_0, x_1, \dots, x_{n-1}\}
\]
donde $x_i$ representa la cantidad de veces que el entero $i$ aparece en la secuencia.

El dominio para cada variable $x_i$ está estrictamente definido en:
\[
D_i = \{0, 1, \dots, n-1\}, \quad \forall i \in \{0, \dots, n-1\}
\]
Se acota superiormente en $n-1$ debido a que ninguna variable puede tomar el valor $n$ (si un número apareciera $n$ veces, ocuparía todas las casillas, obligando a que existan al menos dos números distintos en la secuencia, lo cual es contradictorio).

\paragraph{Restricciones}
\begin{itemize}
    \item \textbf{[R1] Restricción de conteo exacto:} Cada posición $i$ debe coincidir con la sumatoria de coincidencias del valor $i$ en el arreglo:
    \[
    \forall i \in \{0, \dots, n-1\}: \quad x_i = \sum_{j=0}^{n-1} [x_j = i]
    \]
    donde $[\cdot]$ denota el corchete de Iverson.
\end{itemize}

\paragraph{Restricciones redundantes}
El enunciado introduce dos restricciones fundamentales para reducir el árbol de búsqueda:
\begin{itemize}
    \item \textbf{[RR1] Suma total fija:}
    \[
    \sum_{i=0}^{n-1} x_i = n
    \]
    \textit{Justificación:} La secuencia se compone de $n$ posiciones, cada una ocupada por un entero del conjunto $\{0, \dots, n-1\}$. Como cada $x_i$ registra las apariciones del elemento $i$, la suma de todas las frecuencias debe igualar exactamente la cardinalidad de la secuencia.
    
    \item \textbf{[RR2] Suma ponderada diofántica:}
    \[
    (-1)x_0 + 0 \cdot x_1 + 1 \cdot x_2 + \dots + (n-2)x_{n-1} = 0 \iff \sum_{i=0}^{n-1} (i - 1) x_i = 0
    \]
    \textit{Justificación:} Sea $S$ la suma aritmética de todos los elementos presentes en la secuencia. $S$ puede computarse directamente sumando las variables $\sum_{j=0}^{n-1} x_j$, o agrupando por frecuencias como $\sum_{i=0}^{n-1} i \cdot x_i$. Al igualar ambas expresiones:
    \[
    \sum_{i=0}^{n-1} i \cdot x_i = \sum_{j=0}^{n-1} x_j \implies \sum_{i=0}^{n-1} i \cdot x_i - \sum_{i=0}^{n-1} x_i = 0 \implies \sum_{i=0}^{n-1} (i - 1) x_i = 0
    \]
\end{itemize}

\textit{Efecto en la poda del árbol.} Mientras que [R1] propaga de forma tardía (cuando las variables ya están instanciadas), [RR1] y [RR2] son ecuaciones lineales diofánticas. Los propagadores de intervalos asociados (\textit{bounds consistency}) podan las cotas superiores e inferiores de las variables en cada paso de propagación, eliminando ramas inviables sin necesidad de ramificar.

\paragraph{Ruptura de simetrías}
En el problema de la Secuencia Mágica \textbf{no existen simetrías que puedan romperse} (\textbf{[RS1]}). La posición $i$ está rígidamente vinculada a la cuenta de ocurrencias del valor $i$. Intercambiar dos variables $x_i$ y $x_j$ o invertir la secuencia destruiría la relación funcional del modelo. En consecuencia, toda solución es canónica y única en su representación, por lo que no procede aplicar restricciones de simetría.

\subsubsection{Detalles de implementación}
Se estructuraron tres implementaciones incrementales mediante el parámetro \texttt{impl} para cuantificar el impacto de los propagadores lineales:
\begin{enumerate}
    \item \textbf{Implementación 1 (\texttt{impl=1}):} Modelo elemental que únicamente contiene la restricción funcional [R1].
    \item \textbf{Implementación 2 (\texttt{impl=2}):} Incorpora la primera restricción redundante [RR1] junto con [R1].
    \item \textbf{Implementación 3 (\texttt{impl=3}):} Modelo completo que suma ambas restricciones redundantes [R1, RR1, RR2].
\end{enumerate}

\subsubsection{Estrategias de búsqueda}
Se configuraron las siguientes heurísticas de distribución:
\begin{itemize}
    \item \textbf{\texttt{strategy=0}:} Búsqueda estándar de Gecode.
    \item \textbf{\texttt{strategy=1}:} \texttt{input\_order} con \texttt{indomain\_min}.
    \item \textbf{\texttt{strategy=2}:} \texttt{first\_fail} con \texttt{indomain\_min}.
    \item \textbf{\texttt{strategy=3}:} \texttt{dom\_w\_deg} con \texttt{indomain\_min}.
\end{itemize}

\subsubsection{Pruebas}
Los experimentos para la Secuencia Mágica se ejecutaron sobre seis instancias de tamaño progresivo ($n \in \{3, 4, 7, 10, 20, 50\}$). Los resultados obtenidos del archivo \texttt{results/secuencia\_gecode.csv} se sintetizan en la Tabla~\ref{tab:secuencia-resultados}. Además, con la opción \texttt{-a} el modelo enumera todas las secuencias de cada longitud: dos para $n = 4$ ($[1, 2, 1, 0]$ y $[2, 0, 2, 0]$), una para $n = 7$, $10$, $20$ y $50$, y ninguna para $n = 3$. Un recorrido exhaustivo en Python, independiente de MiniZinc, confirmó estos conteos para $n \le 7$.

\begin{table}[htbp]
\small
\caption{Resultados experimentales de la Secuencia Mágica para diferentes órdenes $n$}
\label{tab:secuencia-resultados}
\begin{tabular}{lcccrrr}
\toprule
\textbf{Instancia} & \textbf{Impl.} & \textbf{Estrategia} & \textbf{Estado} & \textbf{Nodos} & \textbf{Fallos} & \textbf{Tiempo (s)} \\
\midrule
secuencia\_01 ($n=4$)  & 1 & 1 (\texttt{input\_order}) & SAT & 7 & 2 & 0,0006 \\
secuencia\_01 ($n=4$)  & 2 & 1 (\texttt{input\_order}) & SAT & 7 & 2 & 0,0001 \\
secuencia\_01 ($n=4$)  & 3 & 2 (\texttt{first\_fail})  & SAT & 3 & 0 & 0,0001 \\
\midrule
secuencia\_03 ($n=10$) & 1 & 1 (\texttt{input\_order}) & SAT & 47 & 22 & 0,0003 \\
secuencia\_03 ($n=10$) & 1 & 2 (\texttt{first\_fail})  & SAT & 93 & 44 & 0,0005 \\
secuencia\_03 ($n=10$) & 2 & 2 (\texttt{first\_fail})  & SAT & 29 & 12 & 0,0002 \\
secuencia\_03 ($n=10$) & 3 & 2 (\texttt{first\_fail})  & SAT & 11 & 5 & 0,0002 \\
\midrule
secuencia\_04 ($n=20$) & 1 & 1 (\texttt{input\_order}) & SAT & 127 & 62 & 0,0030 \\
secuencia\_04 ($n=20$) & 1 & 2 (\texttt{first\_fail})  & SAT & 453 & 224 & 0,0057 \\
secuencia\_04 ($n=20$) & 2 & 2 (\texttt{first\_fail})  & SAT & 69 & 32 & 0,0013 \\
secuencia\_04 ($n=20$) & 3 & 2 (\texttt{first\_fail})  & SAT & 16 & 5 & 0,0002 \\
\midrule
secuencia\_05 ($n=50$) & 1 & 1 (\texttt{input\_order}) & SAT & 367 & 182 & 0,1085 \\
secuencia\_05 ($n=50$) & 1 & 2 (\texttt{first\_fail})  & SAT & 2\,733 & 1\,364 & 0,2716 \\
secuencia\_05 ($n=50$) & 2 & 2 (\texttt{first\_fail})  & SAT & 189 & 92 & 0,0135 \\
secuencia\_05 ($n=50$) & 3 & 2 (\texttt{first\_fail})  & SAT & 31 & 5 & 0,0006 \\
\midrule
secuencia\_06 ($n=3$)  & 1 & 1 (\texttt{input\_order}) & UNSAT & 5 & 3 & 0,0001 \\
secuencia\_06 ($n=3$)  & 2 & 1 (\texttt{input\_order}) & UNSAT & 5 & 3 & 0,0001 \\
secuencia\_06 ($n=3$)  & 3 & 0 (\texttt{default})      & UNSAT & 0 & 1 & 0,0000 \\
\bottomrule
\end{tabular}
\end{table}

\subsubsection{Análisis de los árboles de búsqueda}
\begin{itemize}
    \item \textbf{Poda mediante consistencia de cotas (\textit{bounds consistency}):} Las restricciones [RR1] y [RR2] son formalmente ecuaciones lineales diofánticas. A diferencia de [R1], cuyos propagadores solo actúan cuando las variables están fijadas, las restricciones redundantes permiten que el solver ajuste los límites inferior y superior de los dominios en cada propagación analítica, descartando subárboles completos antes de efectuar divisiones (\textit{split}).
    \item \textbf{Detección inmediata en la raíz para $n=3$:} En la instancia insatisfacible \texttt{secuencia\_06}, el modelo base debió explorar $5$ nodos para probar la inexistencia de soluciones. Al activar ambas restricciones redundantes (\texttt{impl=3}), la contradicción matemática fue detectada directamente en la raíz por propagación algebraica ($0$ nodos explorados y $1$ único fallo).
\end{itemize}

\subsubsection{Análisis de resultados}
\begin{itemize}
    \item \textbf{Efecto drástico de las restricciones redundantes en $n=50$:} En la prueba de estrés (\texttt{secuencia\_05}), el modelo elemental con \texttt{first\_fail} requirió $2\,733$ nodos y $1\,364$ fallos con un tiempo de $0,2716$~s. Al incorporar la sumatoria de cardinalidad (\texttt{impl=2}, [RR1]), el árbol descendió a $189$ nodos y $92$ fallos. Al añadir la sumatoria diofántica ponderada (\texttt{impl=3}, [RR2]), la búsqueda colapsó a apenas $31$ nodos y $5$ fallos en $0,0006$~s (una reducción global superior al 98,8\% en exploración y una aceleración de más de $400$ veces).
\end{itemize}

\subsubsection{Conclusiones del ejercicio}
La formulación de la Secuencia Mágica comprueba de forma empírica la teoría de restricciones redundantes: sin alterar el conjunto de soluciones del problema, transforman un CSP susceptible a una búsqueda combinatoria profunda en un proceso de inferencia casi determinista. La integración conjunta de las ecuaciones de suma total y suma ponderada, junto con la heurística \texttt{first\_fail}, permite resolver instancias de tamaño significativo ($n = 50$) en fracciones de milisegundo.
